% GENERATED FILE. Edit canonical lessons, then run npm run generate:books.
% canonical-source-hash: fnv1a64:8e81e17cc9ef1943
% canonical-lessons: KA-C502-sabhya, KA-C502-vidheya, KA-C502-sistu, KA-C502-gamanisu, KA-C502-spandisu

\chapter{Good Manners}
\label{ch:ka-good-manners}
\hlchaptermodality{502}

\begin{chapteropening}
\textbf{By the end of this chapter:} \emph{I can say polite, obedient, discipline, to take note and to react.}

Complete the last lesson of chapter 502: I can say polite, obedient, discipline, to take note and to react.
\end{chapteropening}

\section[\texorpdfstring{\kn{ಸಭ್ಯ} (sabhya)}{ಸಭ್ಯ (sabhya)}]{\kn{ಸಭ್ಯ} (sabhya) — polite; decent}

\label{lesson:KA-C502-sabhya}



\begin{quote}
Before the new one: say the Kannada for harshly frank, then the Kannada for to mimic mockingly.
\end{quote}

\subsection*{You'll want to know: \kn{ಸಭ್ಯ}}
\textbf{\kn{ಸಭ್ಯ}} — \emph{sabhya} — \textquotedblleft{}polite; decent\textquotedblright{}.

\begin{quote}
\kn{ಅವನು} \kn{ಸಭ್ಯ} \kn{ಹುಡುಗ}, \kn{ಎಲ್ಲರ} \kn{ಜೊತೆ} \kn{ಗೌರವದಿಂದ} \kn{ಮಾತಾಡುತ್ತಾನೆ}.
\end{quote}

\begin{quote}
\emph{avanu sabhya huḍuga, ellara jote gauravadinda mātāḍuttāne.}
\end{quote}

\begin{quote}
He is a polite boy; he speaks respectfully with everyone.
\end{quote}

\kn{ಸಭೆ} is an assembly: \kn{ಸಭ್ಯ}, fit to be in company.

\begin{cousinweb}
\textbf{\kn{ಸಭ್ಯ}} (\emph{sabhya}) is borrowed from Sanskrit \emph{sabhya}, \textquotedblleft{}being in an assembly\textquotedblright{}, and so \textquotedblleft{}polite\textquotedblright{} (CDIAL 13171).
\end{cousinweb}

\begin{grammarlens}[title={What you've built}]
One more word for describing what you lived and felt.
\end{grammarlens}

\subsection*{Your turn}
\begin{itemize}
\raggedright
  \item \emph{Say it:} \emph{sabhya}
  \item \emph{Say it:} \emph{sabhya}, once more
  \item \emph{Say it:} \emph{avanu sabhya huḍuga, ellara jote gauravadinda mātāḍuttāne.} — He is a polite boy; he speaks respectfully with everyone.
  \item \emph{Recall:} say the Kannada for harshly frank, then the Kannada for to mimic mockingly, then say \emph{sabhya} again
\end{itemize}

\begin{tcolorbox}[breakable,colback=teal!4,colframe=teal!35!black,title={Before you move on}]
What does \kn{ಸಭ್ಯ} mean? (\textquotedblleft{}Polite; decent\textquotedblright{}.) Say it once more.
\end{tcolorbox}
\section[\texorpdfstring{\kn{ವಿಧೇಯ} (vidhēya)}{ವಿಧೇಯ (vidhēya)}]{\kn{ವಿಧೇಯ} (vidhēya) — obedient; dutiful}

\label{lesson:KA-C502-vidheya}



\begin{quote}
Before the new one: say the Kannada for to mimic mockingly, then the Kannada for polite.
\end{quote}

\subsection*{You'll want to know: \kn{ವಿಧೇಯ}}
\textbf{\kn{ವಿಧೇಯ}} — \emph{vidhēya} — \textquotedblleft{}obedient; dutiful\textquotedblright{}.

\begin{quote}
\kn{ಆ} \kn{ವಿಧೇಯ} \kn{ಮಗ} \kn{ಅಪ್ಪನ} \kn{ಮಾತು} \kn{ಕೇಳುತ್ತಾನೆ}.
\end{quote}

\begin{quote}
\emph{ā vidhēya maga appana mātu kēḷuttāne.}
\end{quote}

\begin{quote}
That dutiful son listens to his father.
\end{quote}

Formal letters end with \kn{ತಮ್ಮ} \kn{ವಿಧೇಯ}, yours faithfully.

\begin{cousinweb}
\textbf{\kn{ವಿಧೇಯ}} (\emph{vidhēya}) is borrowed from Sanskrit \emph{vidheya}, \textquotedblleft{}to be done, to be governed\textquotedblright{}, and so \textquotedblleft{}docile, obedient\textquotedblright{} (Monier-Williams).
\end{cousinweb}

\begin{grammarlens}[title={What you've built}]
One more word for describing what you lived and felt.
\end{grammarlens}

\subsection*{Your turn}
\begin{itemize}
\raggedright
  \item \emph{Say it:} \emph{vidhēya}
  \item \emph{Say it:} \emph{vidhēya}, once more
  \item \emph{Say it:} \emph{ā vidhēya maga appana mātu kēḷuttāne.} — That dutiful son listens to his father.
  \item \emph{Recall:} say the Kannada for to mimic mockingly, then the Kannada for polite, then say \emph{vidhēya} again
\end{itemize}

\begin{tcolorbox}[breakable,colback=teal!4,colframe=teal!35!black,title={Before you move on}]
What does \kn{ವಿಧೇಯ} mean? (\textquotedblleft{}Obedient; dutiful\textquotedblright{}.) Say it once more.
\end{tcolorbox}
\section[\texorpdfstring{\kn{ಶಿಸ್ತು} (śistu)}{ಶಿಸ್ತು (śistu)}]{\kn{ಶಿಸ್ತು} (śistu) — discipline}

\label{lesson:KA-C502-sistu}



\begin{quote}
Before the new one: say the Kannada for polite, then the Kannada for obedient.
\end{quote}

\subsection*{You'll want to know: \kn{ಶಿಸ್ತು}}
\textbf{\kn{ಶಿಸ್ತು}} — \emph{śistu} — \textquotedblleft{}discipline\textquotedblright{}.

\begin{quote}
\kn{ಶಾಲೆಯಲ್ಲಿ} \kn{ಶಿಸ್ತು} \kn{ತುಂಬಾ} \kn{ಮುಖ್ಯ}.
\end{quote}

\begin{quote}
\emph{śāleyalli śistu tumbā mukhya.}
\end{quote}

\begin{quote}
Discipline is very important at school.
\end{quote}

\kn{ಶಿಸ್ತಾಗಿ}, neatly, in good order.

\begin{cousinweb}
\textbf{\kn{ಶಿಸ್ತು}} (\emph{śistu}) is an everyday Kannada word whose origin is not settled; it is often compared with Sanskrit \emph{śiṣṭa}, \textquotedblleft{}taught, disciplined\textquotedblright{} (CDIAL 12479).
\end{cousinweb}

\begin{grammarlens}[title={What you've built}]
One more word for describing what you lived and felt.
\end{grammarlens}

\subsection*{Your turn}
\begin{itemize}
\raggedright
  \item \emph{Say it:} \emph{śistu}
  \item \emph{Say it:} \emph{śistu}, once more
  \item \emph{Say it:} \emph{śāleyalli śistu tumbā mukhya.} — Discipline is very important at school.
  \item \emph{Recall:} say the Kannada for polite, then the Kannada for obedient, then say \emph{śistu} again
\end{itemize}

\begin{tcolorbox}[breakable,colback=teal!4,colframe=teal!35!black,title={Before you move on}]
What does \kn{ಶಿಸ್ತು} mean? (\textquotedblleft{}Discipline\textquotedblright{}.) Say it once more.
\end{tcolorbox}
\section[\texorpdfstring{\kn{ಗಮನಿಸು} (gamanisu)}{ಗಮನಿಸು (gamanisu)}]{\kn{ಗಮನಿಸು} (gamanisu) — to take note; to notice}

\label{lesson:KA-C502-gamanisu}



\begin{quote}
Before the new one: say the Kannada for obedient, then the Kannada for discipline.
\end{quote}

\subsection*{You'll want to know: \kn{ಗಮನಿಸು}}
\textbf{\kn{ಗಮನಿಸು}} — \emph{gamanisu} — \textquotedblleft{}to take note; to notice\textquotedblright{}.

\begin{quote}
\kn{ನಾನು} \kn{ಅವಳ} \kn{ಹೊಸ} \kn{ಕನ್ನಡಕವನ್ನು} \kn{ಗಮನಿಸಿದೆ}.
\end{quote}

\begin{quote}
\emph{nānu avaḷa hosa kannaḍakavannu gamaniside.}
\end{quote}

\begin{quote}
I noticed her new glasses.
\end{quote}

Past (I …): \kn{ಗಮನಿಸಿದೆ}, \emph{gamaniside}.

\kn{ಗಮನ} is attention: \kn{ಗಮನ} \kn{ಕೊಡು}, to pay attention.

\begin{cousinweb}
\textbf{\kn{ಗಮನಿಸು}} (\emph{gamanisu}) is built on Sanskrit \emph{gamana}, \textquotedblleft{}going\textquotedblright{} (CDIAL 4027), which in Kannada came to mean \textquotedblleft{}attention\textquotedblright{}, with the verb ending -\emph{isu}.
\end{cousinweb}

\begin{grammarlens}[title={What you've built}]
One more word for describing what you lived and felt.
\end{grammarlens}

\subsection*{Your turn}
\begin{itemize}
\raggedright
  \item \emph{Say it:} \emph{gamanisu}
  \item \emph{Say it:} \emph{gamanisu}, once more
  \item \emph{Say it:} \emph{nānu avaḷa hosa kannaḍakavannu gamaniside.} — I noticed her new glasses.
  \item \emph{Recall:} say the Kannada for obedient, then the Kannada for discipline, then say \emph{gamanisu} again
\end{itemize}

\begin{tcolorbox}[breakable,colback=teal!4,colframe=teal!35!black,title={Before you move on}]
What does \kn{ಗಮನಿಸು} mean? (\textquotedblleft{}To take note; to notice\textquotedblright{}.) Say it once more.
\end{tcolorbox}
\section[\texorpdfstring{\kn{ಸ್ಪಂದಿಸು} (spandisu)}{ಸ್ಪಂದಿಸು (spandisu)}]{\kn{ಸ್ಪಂದಿಸು} (spandisu) — to react; to respond}

\label{lesson:KA-C502-spandisu}



\begin{quote}
Before the new one: say the Kannada for discipline, then the Kannada for to take note.
\end{quote}

\subsection*{You'll want to know: \kn{ಸ್ಪಂದಿಸು}}
\textbf{\kn{ಸ್ಪಂದಿಸು}} — \emph{spandisu} — \textquotedblleft{}to react; to respond\textquotedblright{}.

\begin{quote}
\kn{ನನ್ನ} \kn{ಪತ್ರಕ್ಕೆ} \kn{ಕಚೇರಿ} \kn{ಬೇಗ} \kn{ಸ್ಪಂದಿಸಿತು}.
\end{quote}

\begin{quote}
\emph{nanna patrakke kacēri bēga spandisitu.}
\end{quote}

\begin{quote}
The office responded quickly to my letter.
\end{quote}

Past (I …): \kn{ಸ್ಪಂದಿಸಿದೆ}, \emph{spandiside}.

To what: -\kn{ಗೆ}. Literally to throb in answer.

\begin{cousinweb}
\textbf{\kn{ಸ್ಪಂದಿಸು}} (\emph{spandisu}) is built on Sanskrit \emph{spanda}, \textquotedblleft{}a tremor, throbbing\textquotedblright{} (CDIAL 13805), with the verb ending -\emph{isu}.
\end{cousinweb}

\begin{grammarlens}[title={What you've built}]
One more word for describing what you lived and felt.
\end{grammarlens}

\subsection*{Your turn}
\begin{itemize}
\raggedright
  \item \emph{Say it:} \emph{spandisu}
  \item \emph{Say it:} \emph{spandisu}, once more
  \item \emph{Say it:} \emph{nanna patrakke kacēri bēga spandisitu.} — The office responded quickly to my letter.
  \item \emph{Recall:} say the Kannada for discipline, then the Kannada for to take note, then say \emph{spandisu} again
\end{itemize}

\begin{tcolorbox}[breakable,colback=teal!4,colframe=teal!35!black,title={Before you move on}]
What does \kn{ಸ್ಪಂದಿಸು} mean? (\textquotedblleft{}To react; to respond\textquotedblright{}.) Say it once more.
\end{tcolorbox}
